% =============================================================
% macros.tex
% 教材用マクロ（役者ボックス・確認問題・比較表・図版・QR）
% 既存マクロのシグネチャは全て維持し、ビジュアルだけ再設計。
% =============================================================

% -------------------------------------------------------------
% 1. サポートサイト URL と QR（一元管理）
% -------------------------------------------------------------
\edef\supportbaseurl{\detokenize{https://morisanyutakun-png.github.io/ITpass_book/}}
\NewDocumentCommand{\supporturl}{m}{\supportbaseurl\detokenize{#1}.html}

% -------------------------------------------------------------
% 2. 役者ボックスの共通定義
% 「左バー＋小さなラベルタブ＋淡色背景」で統一する。
% 共通の見た目を actorbox として定義し、色・ラベルだけ差し替える。
% -------------------------------------------------------------
% #1: ラベル, #2: ラベル背景色, #3: 本文背景色, #4: バー色, #5: 文字色
\newtcolorbox{actorbox}[5]{%
  enhanced, breakable,
  colback=#3,
  colframe=#3,          % 外枠は背景と同色（見えないように）
  boxrule=0pt,
  leftrule=\actorBarWidth,
  colbacktitle=#2,
  coltitle=white,
  fonttitle=\sffamily\footnotesize\bfseries,
  title=#1,
  attach boxed title to top left={xshift=-\actorBarWidth, yshift=-6pt},
  boxed title style={
    colframe=#2, boxrule=0pt,
    sharp corners, arc=0pt,
    left=\labelPadX, right=\labelPadX,
    top=\labelPadY, bottom=\labelPadY,
  },
  overlay unbroken and first={
    \draw[#4, line width=\actorBarWidth] (frame.north west) -- (frame.south west);
  },
  arc=\actorArc, outer arc=0pt,
  sharp corners=northwest,
  left=12pt, right=10pt, top=12pt, bottom=8pt,
  fontupper=\color{#5}\small,
  before skip=10pt, after skip=10pt,
}

% -------------------------------------------------------------
% 3. 定義ボックス（要点）
% -------------------------------------------------------------
\newtcolorbox{definitionbox}[1][要点]{%
  enhanced, breakable,
  colback=pointBG, colframe=pointBG, boxrule=0pt,
  leftrule=\actorBarWidth,
  colbacktitle=pointBar, coltitle=white,
  fonttitle=\sffamily\footnotesize\bfseries,
  title=\,#1\,,
  attach boxed title to top left={xshift=-\actorBarWidth, yshift=-6pt},
  boxed title style={
    colframe=pointBar, boxrule=0pt,
    sharp corners, arc=0pt,
    left=\labelPadX, right=\labelPadX,
    top=\labelPadY, bottom=\labelPadY,
  },
  overlay unbroken and first={
    \draw[pointBar, line width=\actorBarWidth]
      (frame.north west) -- (frame.south west);
  },
  arc=\actorArc, outer arc=0pt,
  left=12pt, right=10pt, top=14pt, bottom=8pt,
  fontupper=\color{themeInk}\small,
  before skip=10pt, after skip=10pt,
}
% 互換：\definitionbx{タイトル}{本文}
\newcommand{\definitionbx}[2]{\begin{definitionbox}[#1]#2\end{definitionbox}}

% -------------------------------------------------------------
% 4. つまずきポイント
% -------------------------------------------------------------
\newtcolorbox{pitfallbox}[1][つまずき]{%
  enhanced, breakable,
  colback=pitBG, colframe=pitBG, boxrule=0pt,
  leftrule=\actorBarWidth,
  colbacktitle=pitBar, coltitle=white,
  fonttitle=\sffamily\footnotesize\bfseries,
  title=\,#1\,,
  attach boxed title to top left={xshift=-\actorBarWidth, yshift=-6pt},
  boxed title style={
    colframe=pitBar, boxrule=0pt,
    sharp corners, arc=0pt,
    left=\labelPadX, right=\labelPadX,
    top=\labelPadY, bottom=\labelPadY,
  },
  overlay unbroken and first={
    \draw[pitBar, line width=\actorBarWidth]
      (frame.north west) -- (frame.south west);
  },
  arc=\actorArc, outer arc=0pt,
  left=12pt, right=10pt, top=14pt, bottom=8pt,
  fontupper=\color{themeInk}\small,
  before skip=10pt, after skip=10pt,
}
\newcommand{\pitfall}[2]{\begin{pitfallbox}[#1]#2\end{pitfallbox}}

% -------------------------------------------------------------
% 5. 比較ボックス
% -------------------------------------------------------------
\newtcolorbox{comparisonbox}[1][比較]{%
  enhanced, breakable,
  colback=cmpBG, colframe=cmpBG, boxrule=0pt,
  leftrule=\actorBarWidth,
  colbacktitle=cmpBar, coltitle=white,
  fonttitle=\sffamily\footnotesize\bfseries,
  title=\,#1\,,
  attach boxed title to top left={xshift=-\actorBarWidth, yshift=-6pt},
  boxed title style={
    colframe=cmpBar, boxrule=0pt,
    sharp corners, arc=0pt,
    left=\labelPadX, right=\labelPadX,
    top=\labelPadY, bottom=\labelPadY,
  },
  overlay unbroken and first={
    \draw[cmpBar, line width=\actorBarWidth]
      (frame.north west) -- (frame.south west);
  },
  arc=\actorArc, outer arc=0pt,
  left=10pt, right=8pt, top=14pt, bottom=6pt,
  fontupper=\color{themeInk}\small,
  before skip=10pt, after skip=10pt,
}
\newcommand{\comparisonbx}[2]{\begin{comparisonbox}[#1]#2\end{comparisonbox}}

% -------------------------------------------------------------
% 6. 1 枚まとめ
% -------------------------------------------------------------
\newtcolorbox{summarybox}[1][この節のまとめ]{%
  enhanced, breakable,
  colback=sumBG, colframe=sumBG, boxrule=0pt,
  leftrule=\actorBarWidth,
  colbacktitle=sumBar, coltitle=white,
  fonttitle=\sffamily\footnotesize\bfseries,
  title=\,#1\,,
  attach boxed title to top left={xshift=-\actorBarWidth, yshift=-6pt},
  boxed title style={
    colframe=sumBar, boxrule=0pt,
    sharp corners, arc=0pt,
    left=\labelPadX, right=\labelPadX,
    top=\labelPadY, bottom=\labelPadY,
  },
  overlay unbroken and first={
    \draw[sumBar, line width=\actorBarWidth]
      (frame.north west) -- (frame.south west);
  },
  arc=\actorArc, outer arc=0pt,
  left=12pt, right=10pt, top=14pt, bottom=10pt,
  fontupper=\color{themeInk}\small,
  before skip=12pt, after skip=12pt,
}
\newcommand{\onelinesummary}[2][この節のまとめ]{\begin{summarybox}[#1]#2\end{summarybox}}

% -------------------------------------------------------------
% 7. 学習アドバイス（\advice{本文}）
%   本試験参照用の adviceinline とは別スタイルにして混同を避ける
% -------------------------------------------------------------
\newtcolorbox{studyadvice}{%
  enhanced, breakable,
  colback=adviceBG, colframe=adviceBG, boxrule=0pt,
  leftrule=\actorBarWidth,
  colbacktitle=adviceBar, coltitle=white,
  fonttitle=\sffamily\footnotesize\bfseries,
  title=\,学習アドバイス\,,
  attach boxed title to top left={xshift=-\actorBarWidth, yshift=-6pt},
  boxed title style={
    colframe=adviceBar, boxrule=0pt,
    sharp corners, arc=0pt,
    left=\labelPadX, right=\labelPadX,
    top=\labelPadY, bottom=\labelPadY,
  },
  overlay unbroken and first={
    \draw[adviceBar, line width=\actorBarWidth]
      (frame.north west) -- (frame.south west);
  },
  arc=\actorArc, outer arc=0pt,
  left=12pt, right=10pt, top=14pt, bottom=8pt,
  fontupper=\color{themeInk}\small,
  before skip=10pt, after skip=10pt,
}
\newcommand{\advice}[1]{\begin{studyadvice}#1\end{studyadvice}}

% -------------------------------------------------------------
% 8. 本試験の出題パターン（\begin{adviceinline} 用）
%   章中で \begin{adviceinline}...\end{adviceinline} の raw 利用あり。
%   紫紺アクセントで「学習アドバイス」と視覚的に区別する。
% -------------------------------------------------------------
\newtcolorbox{adviceinline}{%
  enhanced, breakable,
  colback=examBG, colframe=examBG, boxrule=0pt,
  leftrule=\actorBarWidth,
  colbacktitle=examBar, coltitle=white,
  fonttitle=\sffamily\footnotesize\bfseries,
  title=\,本試験\,,
  attach boxed title to top left={xshift=-\actorBarWidth, yshift=-6pt},
  boxed title style={
    colframe=examBar, boxrule=0pt,
    sharp corners, arc=0pt,
    left=\labelPadX, right=\labelPadX,
    top=\labelPadY, bottom=\labelPadY,
  },
  overlay unbroken and first={
    \draw[examBar, line width=\actorBarWidth]
      (frame.north west) -- (frame.south west);
  },
  arc=\actorArc, outer arc=0pt,
  left=12pt, right=10pt, top=14pt, bottom=8pt,
  fontupper=\color{themeInk}\small,
  before skip=10pt, after skip=10pt,
}

% -------------------------------------------------------------
% 9. 確認問題（演習コーナーとして再設計）
%   \checkquestion{Q}{A}{B}{C}{D}{正解}{解説} の互換シグネチャ。
%   「演習バー」→ 問題 → 選択肢 → 解答カード の 4 層に分離。
% -------------------------------------------------------------
\newtcolorbox{answerbox}{%
  enhanced, breakable,
  colback=ansBG, colframe=ansBG, boxrule=0pt,
  leftrule=\actorBarWidth,
  colbacktitle=ansBar, coltitle=white,
  fonttitle=\sffamily\footnotesize\bfseries,
  title=\,解答・解説\,,
  attach boxed title to top left={xshift=-\actorBarWidth, yshift=-6pt},
  boxed title style={
    colframe=ansBar, boxrule=0pt,
    sharp corners, arc=0pt,
    left=\labelPadX, right=\labelPadX,
    top=\labelPadY, bottom=\labelPadY,
  },
  overlay unbroken and first={
    \draw[ansBar, line width=\actorBarWidth]
      (frame.north west) -- (frame.south west);
  },
  arc=\actorArc, outer arc=0pt,
  left=12pt, right=10pt, top=14pt, bottom=8pt,
  fontupper=\color{themeInk}\small,
  before skip=4pt, after skip=12pt,
}

% 「演習」バー（section-like ミニ見出し）
\newcommand{\quizhead}{%
  \par\addvspace{12pt}%
  \noindent
  \tikz[baseline=(q.base)]{
    \node[
      fill=quizBar, text=white,
      inner xsep=8pt, inner ysep=2.5pt,
      font=\sffamily\footnotesize\bfseries, rounded corners=1pt
    ] (q) {\,演習\,};
  }%
  \hspace{6pt}%
  {\color{themeHair}\leaders\hrule height 0.4pt\hfill\null}%
  \par\vspace{4pt}%
}

\NewDocumentCommand{\checkquestion}{m m m m m m m}{%
  \begingroup
  \quizhead
  \noindent\textbf{\color{themeMain}Q.} #1\par
  \vspace{4pt}
  \begin{enumerate}[label=\textbf{\color{themeSub}\Alph*.},
                   leftmargin=2.2em, itemsep=2pt, topsep=3pt]
    \item #2
    \item #3
    \item #4
    \item #5
  \end{enumerate}
  \vspace{2pt}
  \begin{answerbox}
    \textbf{\color{ansBar}正解}\quad #6 \par
    \vspace{2pt}
    \textbf{\color{ansBar}解説}\quad #7
  \end{answerbox}
  \endgroup
}

% -------------------------------------------------------------
% 10. 章末まとめ見出し（既存互換）
% -------------------------------------------------------------
\newcommand{\chaptersummary}[1]{%
  \section*{この章のまとめ}%
  \addcontentsline{toc}{section}{この章のまとめ}%
  #1
}

% -------------------------------------------------------------
% 11. 章扉のリード文（章の"役割"を示す 1 段落）
%   \chapter{...} の直後に任意で置く。
% -------------------------------------------------------------
\newtcolorbox{chapterleadbox}{%
  enhanced, breakable,
  colback=chapLeadBG, colframe=chapLeadBG, boxrule=0pt,
  leftrule=\actorBarWidth,
  overlay unbroken and first={
    \draw[chapTagLine, line width=\actorBarWidth]
      (frame.north west) -- (frame.south west);
  },
  arc=\actorArc, outer arc=0pt,
  left=14pt, right=12pt, top=10pt, bottom=10pt,
  fontupper=\color{themeInk}\small,
  before skip=0pt, after skip=16pt,
}
\newcommand{\chapterlead}[1]{\begin{chapterleadbox}#1\end{chapterleadbox}}

% -------------------------------------------------------------
% 12. 図版マクロ（キャプション＋ラベルタグ）
% -------------------------------------------------------------
\NewDocumentCommand{\figurebox}{m m m}{%
  \begin{figure}[ht]
    \centering
    {\color{themeHair}\rule{0.92\linewidth}{0.3pt}}\par
    \vspace{4pt}
    \adjustbox{max width=0.94\linewidth}{#3}\par
    \vspace{3pt}
    \noindent\hspace*{4\%}%
    \parbox{0.92\linewidth}{%
      \small\color{themeMuted}%
      \tikz[baseline=(c.base)]{
        \node[fill=themeAccent, text=white,
              inner xsep=5pt, inner ysep=1.5pt,
              font=\sffamily\scriptsize\bfseries,
              rounded corners=1pt] (c) {\,図\,\thechapter.\arabic{figure}\,};
      }%
      \ #1\label{#2}%
    }\par
    \vspace{2pt}
    {\color{themeHair}\rule{0.92\linewidth}{0.3pt}}
  \end{figure}%
}

% -------------------------------------------------------------
% 13. 比較表テンプレ（ヘッダ色付き・罫線弱め）
% -------------------------------------------------------------
\NewDocumentCommand{\comptable}{m m}{%
  \begin{center}
  \small
  \renewcommand{\arraystretch}{1.2}%
  \arrayrulecolor{themeHair}%
  \begin{tabularx}{0.96\linewidth}{@{}>{\color{themeSub}\bfseries\sffamily\footnotesize}l X X@{}}
    \arrayrulecolor{cmpBar}\toprule
    #1 \\
    \arrayrulecolor{themeHair}\midrule
    #2 \\
    \arrayrulecolor{cmpBar}\bottomrule
  \end{tabularx}
  \arrayrulecolor{black}%
  \end{center}%
}

% -------------------------------------------------------------
% 14. 強調タグ（本文中の小さな強調）
% -------------------------------------------------------------
\newcommand{\termtag}[1]{\textcolor{themeAccent}{\textbf{#1}}}
\newcommand{\altterm}[1]{\textcolor{themeSub}{\textsf{#1}}}

% -------------------------------------------------------------
% 15. QR サポートカード（枠を細く・背景なしで"しおり"風に）
% -------------------------------------------------------------
\makeatletter
\newcommand{\itp@qrurl}{}
\NewDocumentCommand{\supportqr}{m m}{%
  \begingroup
  \edef\itp@qrurl{\supportbaseurl\detokenize{#1}.html}%
  \par\addvspace{10pt}%
  \noindent
  \begin{minipage}{\linewidth}
    % 上部：細い罫線＋ラベル
    {\color{themeAccent}\rule[0.4em]{3pt}{10pt}}%
    \hspace{4pt}%
    {\sffamily\footnotesize\bfseries\color{themeAccent}SUPPORT}%
    \hspace{6pt}%
    {\small\color{themeSub}#2}%
    \hfill
    {\scriptsize\color{themeMuted}サポートサイト}%
    \par\vspace{3pt}%
    {\color{themeHair}\rule{\linewidth}{0.3pt}}\par
    \vspace{4pt}
    % 本体：QR 左・説明右
    \noindent
    \begin{minipage}[c]{26mm}
      \centering\qrcode[height=22mm]{\itp@qrurl}
    \end{minipage}%
    \hfill
    \begin{minipage}[c]{\dimexpr\linewidth-32mm\relax}
      \footnotesize
      スマートフォンのカメラで読み取ると、本書のサポートサイト該当ページへ移動します。補足図解・更新履歴・確認テストなどを配信しています。\par
      \vspace{3pt}
      {\tiny\urlstyle{tt}\color{themeMuted}\expandafter\nolinkurl\expandafter{\itp@qrurl}}
    \end{minipage}%
    \par\vspace{4pt}%
    {\color{themeHair}\rule{\linewidth}{0.3pt}}%
  \end{minipage}%
  \par\addvspace{10pt}%
  \endgroup
}
\makeatother
